% ============================================================================
% PLATZHALTER -- Bild 6: Heatmap |e^{st} F(s)| mit Bromwich- und
%                        Talbot-Pfad
% ============================================================================
% REGIEANWEISUNG FUER DIE GRAFIK-TOOLCHAIN (matplotlib empfohlen,
% Export als PDF; Schrift Times 10pt; KEINE transparenten PNG-Teile):
%
% DATENGRUNDLAGE (F(s)=1/s, t=1):
% - Gitter der s-Ebene: Re(s) in [-15, 5], Im(s) in [-12, 12],
%   Aufloesung >= 800x800.
% - Farbwert: log10( |e^{s*t} / s| ), geclippt auf z.B. [-8, 2].
%
% DARSTELLUNG:
% - 2D-CONTOUR-HEATMAP (pcolormesh oder contourf mit vielen Stufen),
%   sequentielle Colormap (z.B. viridis); COLORBAR rechts mit
%   Beschriftung "$\log_{10}|e^{st}F(s)|$".
% - Zusaetzlich duenne weisse Konturlinien bei ganzzahligen log10-Stufen,
%   damit die "Landschaft" (Bergkamm rechts, Tal links) plastisch wird.
% - SINGULARITAET s=0: kleines weisses/rotes Kreuz mit Label "Pol".
% - BROMWICH-PFAD: vertikale ROTE Linie bei Re(s)=1 ueber die volle
%   Bildhoehe, Pfeilspitze nach oben, Label "$\Gamma_B$" (rot).
% - TALBOT-PFAD: BLAUE Kurve s(theta)=sigma+lambda(theta*cot(theta)
%   + i*beta*theta) mit denselben Parametern wie in Bild 5 (sigma=0,
%   lambda=2, beta=1; Formparameter beta, nicht nu beschriften),
%   Pfeilspitzen in Durchlaufrichtung (von unten nach
%   oben), Label "$\Gamma_T$" (blau).
%   Die Kurve muss erkennbar ins dunkle "Tal" links abtauchen und den
%   Pol rechts umschliessen.
% - Textannotation "Tal: $|e^{st}|$ exponentiell klein" im linken
%   Bildbereich (weiss, small); optional "Bergkamm" rechts.
% - Achsenbeschriftung "Re(s)", "Im(s)"; kein Bildtitel.
%
% ALTERNATIVE (falls 3D gewuenscht): 3D-Surface derselben Daten mit
% Blick von schraeg oben, Pfade auf die Flaeche projiziert; 2D-Version
% ist fuer den Druck aber besser lesbar und wird bevorzugt.
%
% GROESSE: ca. 11cm x 8cm inkl. Colorbar, 300dpi falls Rasterexport.
% ============================================================================
\scalebox{0.85}{
\begin{tikzpicture}

% ==============================
% 1. FARBEN DEFINIEREN
% ==============================
\definecolor{plotblue}{RGB}{0, 0, 200}
\definecolor{plotred}{RGB}{220, 30, 30}
\definecolor{plotgray}{RGB}{140, 140, 140}
\definecolor{plotbg}{RGB}{244, 244, 252}

% PGFPLOTS SETUP
\pgfplotsset{table/search path={.,Latex/images,images,papers/laplace/images}}

\begin{axis}[
    width=11cm,
    height=9.5cm,
    enlargelimits=false,
    axis on top,
    xlabel={$\operatorname{Re}(s)$},
    ylabel={$\operatorname{Im}(s)$},
    xmin=-15, xmax=5,
    ymin=-12, ymax=12,
    tick label style={font=\small},
    label style={font=\normalsize},
    colorbar,
    colormap/viridis,
    point meta min=-8,
    point meta max=2,
    colorbar style={
        ylabel={$\log_{10}|e^{st}F(s)|$},
        ytick={-8,-6,-4,-2,0,2},
        ylabel style={font=\normalsize},
        tick label style={font=\small}
    }
]

% Hintergrund-Heatmap (Bild, darauf projiziert)
\addplot graphics [xmin=-15, xmax=5, ymin=-12, ymax=12] {papers/laplace/images/heatmap_bg.png};

% Singularitaet s=0 (Pol)
\draw[thick, white, line cap=round] (axis cs:-0.4,-0.4) -- (axis cs:0.4,0.4) (axis cs:-0.4,0.4) -- (axis cs:0.4,-0.4);
\draw[thick, plotred, line cap=round] (axis cs:-0.2,-0.2) -- (axis cs:0.2,0.2) (axis cs:-0.2,0.2) -- (axis cs:0.2,-0.2);
\node[white, left, font=\small, xshift=-2pt] at (axis cs:-0.2, 0) {\textbf{Pol}};

% Bromwich-Gerade (rot)
\draw[thick, plotred, decoration={markings, mark=at position 0.65 with {\arrow{Stealth[length=3mm, width=2mm]}}}, postaction={decorate}] 
    (axis cs:1, -12) -- (axis cs:1, 12);
\node[plotred, right, font=\large] at (axis cs:1.2, 9) {$\Gamma_B$};

% Talbot-Kontur (blau)
\addplot[
    thick, plotblue,
    decoration={markings, mark=at position 0.35 with {\arrow{Stealth[length=3mm, width=2mm]}}},
    decoration={markings, mark=at position 0.65 with {\arrow{Stealth[length=3mm, width=2mm]}}},
    postaction={decorate}
    ] table[col sep=comma, x=re, y=im] {papers/laplace/images/talbot_contour_heatmap.csv};
\node[plotblue, left, font=\large] at (axis cs:-3, 6) {$\Gamma_T$};

% Annotation Tal removed

% ==============================
% 2. LEGENDE
% ==============================
\filldraw[fill=white, draw=black, thick, rounded corners=3pt] (axis cs:-14.5, 11.5) rectangle (axis cs:-2.5, 8.5);
\draw[thick, plotred, decoration={markings, mark=at position 0.55 with {\arrow{Stealth[length=3mm, width=2mm]}}}, postaction={decorate}] (axis cs:-14.0, 10.5) -- (axis cs:-12.5, 10.5);
\node[right] at (axis cs:-12.0, 10.5) {Bromwich-Gerade $\Gamma_B$};
\draw[thick, plotblue, decoration={markings, mark=at position 0.55 with {\arrow{Stealth[length=3mm, width=2mm]}}}, postaction={decorate}] (axis cs:-14.0, 9.5) -- (axis cs:-12.5, 9.5);
\node[right] at (axis cs:-12.0, 9.5) {Talbot-Kontur $\Gamma_T$};

\end{axis}
\end{tikzpicture}
}
